\documentclass[12pt]{article}

\usepackage[a4paper, total={6in, 8in}]{geometry}
\usepackage{textcomp}

\begin{document}
\setlength{\parindent}{0pt}

\def \t {\textrightarrow}
\def \v {\vspace{1em}}
\def \bi {\begin{itemize}}
\def \ei {\end{itemize}}
\def \s[#1] {\section*{#1}}
\def \ss[#1] {\subsection*{#1}}
\def \sss[#1] {\subsubsection*{#1}}

\s[Croce]
Neoidealista, si avvicina a Hegel attraverso Marx \\
Dimora da Spaventa a Roma, e qua è un ambiente ricco intellettualmente \t incontra il mondo culturale italiano \\
Pensa anche che Hegel sia troppo complicato e che quindi sia da allontanare \\
Poi torna a Napoli \t mentre a Roma le discussioni erano molto politiche, a Napoli l'ambiente era + erudito \\
Studia e viaggia, ma non consegue titolo accademico \t l'unica scintilla è quella del marxismo \\
Inizia a condividere questo pensiero, ma per poco tempo sara vicino a queste posizioni \t scopre anche un Hegel diverso da quello che aveva studiato all'inizio \\
Egli trova in Marx dei punti deboli, e vuole risalire alle fonti \t arriva a un Hegel + teoretico, non quello che aveva conosciuto a Roma \\
Ripensa la filosofia di hegel, e nel 1906 pubblica "Ciò che è vivo e ciò che è morto nella filosofia di Hegel" \t nel frattempo aveva gia scritto "L'estetica", opera + famosa del 1902 \\
Nell'estate del 1903 Croce aveva anche creato una rivista, insieme a Gentile \t con il quale era amico, finche Gentile non aderisce al fascismo \\
Nel 1910 Croce diventa senatore e ministro dell'istruzione tra il 20 e il 21 \t inzia una riforma scolastica che poi finisce Gentile durante il periodo fascista \\
Dopo delitto Matteoti, Croce assume posizione antifascista \t e sarà il primo firmatario del manifesto dei letterati antifascisti \\
Dopo fascismo diventa membro dell'assemblea costituente \t nel ? fonda istituto per gli studi storici, e muore nel 52 \t lavora molto per la casa editrice Laterza

\v

Scrive opere teoretiche, saggi filosofici + generici, scrive 54 volumi di storia letteraria e politica + 12 volumi di scritti vari \t scrive molto

\ss[Ciò che è vivo e ciò che è morto nella filosofia di Hegel]
Manifesto dell'idealismo di Croce \t dice che Hegel ha capito molte cose, come la vera dimensione del pensiero filosofico, che il concetto è il nome + alto dello spirito, dialettica etc. \\
Però non ha capito che non tutta la realtà è sintesi di opposti \t non tutto nella realtà puo essere messa in relazione dialettica \\
Il vero e il male non sono antitetici, sono due cose diverse \t la dialettica non è solo sintesi di opposti, ma anche nesso di distinti = all'interno della dialettica, ci sono degli ambiti che si danno sempre tutti insieme e sono dei momenti dello spirito che non sono sintetizzabili \\
Questi ambiti sono distinti \t la dialettica è anche nesso di distinti, legame tra di loro \t all'interno della vita si vedono questi 4 momenti ? \\
Questo è l'errore dell hegelismo \t la nuova dialettica è molto diversa quindi, e nasce da nuove considerazioni: \\
Lo spirito è attraversato da:
\bi
  \item attivita teoretica (conoscitiva)
  \item attivita pratica (etica, legata all'azione)
\ei
Le due attivita si possono orientare verso oggetti diversi \t l'oggetto del teoretico e del pratico può essere sia particolare sia universale \t nascono quindi 4 distinti, 4 momenti della vita dello spirito che si danno sempre insieme e non sono sintetizzabili
\bi
  \item teoretica rivolta all'individuale = estetica, coppia di opposti bello/brutto, tra cui avviene la sintesi
  \item teoretica rivolta all'universlae = logica, coppia vero/falso
  \item pratica all'individuale = economia, utile/dannoso (nel senso dell'utilità)
  \item pratica all'universale = etica, opposti sono bene/male
\ei
Vita spirito è fatto da questi 4 distinti all'interno dei quali avviene la sintesi degli opposti \t la nostra indagine quindi deve seguire questi ambiti e fare ricerca settoriale in questi, se si vuole conoscere lo spirito

\ss[L'estetica + Breviario di estetica]
Primo ambito che analizza \t ed è il suo capolavoro quest'opera \\
La domanda con cui si apre è cosa sia l'arte \t tutti gli uomini l'hanno incontrata nella loro vita, è la cosa che tutti sanno si puo rispondere \\
La filosofia ha invece il compito di portare a comprensione chiara cio che tutti conoscono ma non sanno \t analizza le domande che sembrano banali ma non lo sono \\
L'arte dice che è conoscenza intuitiva \t si è pensato che l'intuizione fosse qualcosa di ceco e che l'intelletto dovesse sostenerla, ma l'intuizione è autonoma \t e vede le cose chiaramente quanto l'intelletto \\
L'arte è conoscenza intuitiva = l'artista non fa ragionamento che trasforma nell'opera, ma intuisce quello che immediatamente trasferisce nella sua opera \\
Nell'intuzione si mettono insieme tutte le nostre conoscenze, e l'artista reinterpreta questo bagaglio nella sua scintilla intuitiva che gli permette di realizzare l'opera

\v

(2) "L'intutuizione estetica è sempre anche un'espressione estetica" \t l'espressione e intuizione sono in realtà stessa cosa, il modo in cui intuisco è il modo in cui realizzo \\
Gli artisti però sono uomini \t l'intuizione estetica appartiene a tutti gli uomini, altrimenti non si sarebbe in grado di cogliere bellezza di un opera \t il genio pero ha maggiore quantita intuitiva \t differenza quantitativa, e non qualitativa \\
Questo spiega perche l'artista è umano, perchè gli uomini comprendo l'arte \\
Quindi, io esprimo anche tanto quanto intuisco \t se creo un dipinto, poesia etc., la mia intuizione estetica corrisponde alla mia espressione

\v

(3) L'intuzione puo solo nascere da un sentimento, non è il concetto che fa nascere intuizione artistica \t è il sentimento che conferisce la leggerezza del simbolo all'arte \\
Io colgo l'opera d'arte ma colgo anche qualcosa che va oltre \\
Il sentimento è lirico, l'intuizione è lirica (non è un attributo dell'estetica, la liricità è sinonimo di intuizione)

\v

Cosa però definiamo veramente opera d'arte? \t l'opera d'arte è tale perchè è una sintesi perfetta di forma e contenuto \t è arte per il contenuto e per la sua forma \\
Un contenuto non formato e una forma senza contenuto non sono opere d'arte \\
L'arte ha un carattere di universalità \t quello che esprime va al di la del dettaglio particolare che la poesia sta descrivendo \\
Poi croce afferma che i generi letterari non esistono \t l'arte è sempre unica in tutte le sue manifestazioni, ci sono solo degli schemi comodi all'intelletto per classificare, che però non esistono nell'arte \\
Non c'è rilevanza dal punto di vista estetico \t se si considerano di valenza estetica la classificazione, si cade nell'iperrazionalismo \\
Anche le tecniche artistiche non appartengono \t corrispondono a come comunico l'espressione \t le uso semplicemente per rendere visibile la mia espressione \\
Il linguaggio è essenzialmente espressione, che è anche espressione dell'arte \t quindi il linguaggio è una creazione estetica \t il linguaggio come espressione libera, la forma del parlare

\ss[Logica]
Qua è molto hegeliano, introduce l'"univesale concreto" \\
La logica è la scienza del concetto puro \t che esprime l'unica essenza della realta \\
Il concetto è unico ed è il nome + perfetto dell'assoluto \t ma anche della realtà \t il concetto vale sul piano logico, ma ha anche dimensione ontologica \t è il nome che do allo spirito, ma è anche lo spirito \t per questo il concetto è l'universale concreto \\
Tutti i concetti universali sono pseudoconcetti \t quando dico cane, sto facendo una riflessione di carattere economico: mi serve usare i concetti perchè mi devo muovere nel mondo \\
Tutte le scienze empiriche e matematiche, non hanno valore teoretico \t e addirittura stanno nell'attivià pratica, non descrivono vera essenza dellà realtà \\
Per l'idealismo dell ottocento questi pseudoconcetti erano opera dell'intelletto \t Croce invece afferma che non nascono da attivià teoretica, l'intelletto coincide con la nostra razionalita che sa che esiste un unico concetto \t uomo usa dei nomi solo per utilià \\
Non sono elaborazioni teoretiche, ma pratiche \\
Nella logica pone una delle sue tesi + tipiche \t che porta a parlare dello storicismo della sua filosofia \t il giudizio individuale e quello definitorio coincidono \\
Anche il giudizio cambia di senso in questra prospettiva:
\bi
  \item il giudizio definitorio è quello tipico della filosofia, es. "l'arte è intuizione lirica" = universale che da definizione 
  \item il giudizio individuale è tipico della storia \t "l'orlando furioso è un opera d'arte"
\ei
Per dire questo fatto particolare devo pero aver prima definito l'universale \t devo prima definire opera d'arte \\
Tra filosofia e storia coincidono quindi \t e si puo dire che la filosofia sta nella storia: per definire il particolare, devo prima definire l'universale \\
Quindi l'unica conoscenza del reale è la storia \t è una sintesi a priori tra l'universale e il particolare, tra il fatto e la categoria che definisce il fatto, che devo aver preventivamente definito \\
La storia è l'insieme dell'evento e la sua valutazione filosofica \\
Quindi se io studio la storia ho la massima conoscenza del reale \t perchè conoscere la storia significa conoscere lo spirito \t la storia è anche la visione perfetta della razionalità \\
È presente una necessità della storia, che è data dalla natrua dello spirito che è la sostanza della storia stessa

\v

Se così è allora il se storico è ridicolo \t se storia è attraversata da questa razionalità, allora non poteva andare in altro modo la storia \t storia esprime razionalità necessaria \\
Allora anche il fascismo però è un evento della storia \t che devo accettare come una manifestazione dello spirito \t anche se poi lui lo combattera \\
Gentile invece è + corente teoreticamente

\v

Con croce si ha storicismo assoluto \t e dice che non esiste mai una storia passata, il periodo al quale si riferisce è sempre attuale \t la storia dello spirito è sempre quella, lo spirito diviene portando dietro il suo passato \\
Non ci sono epoche storiche distinte, non esiste un periodo che non ha influenza sul presente \t è una storia che continua nel presente, e il soggetto di questa storia è lo spirito \\
Anche noi ci portiamo dietro la storia, perche noi siamo manifestazione dello spirito \t noi siamo il risultato di tutta quella storia \t molto diverso da Bergson \\
Qua siamo risultato di storia necessaria, e non contiamo niente \t mentre Bergson dice che ci portiamo dietro la nostra storia, che non è necessaria

\ss[Gentile]
Gentile fa riforma dell'hegelismo che è chiamata l'attualismo \t è quella forma di hegelismo che dice che lo spirito è sempre in atto \t lo spirito è quell atto che pone il suo oggetto nell'altro \\
Lo spirito si pone sempre in atto, perche questo ponersi genera molteplicita di oggetti che vengono colti altro da se, ma che vengono riassorbiti con la sintesi \\
L'attualismo è una semplificazione dell'hegelismo, che sfronda hegel riptendolo \t lo spirito è sempre in atto, e in quanto in atto oppone sempre a se un oggetto che poi va a recuperare \\
Gentile nl 22 diventa senatore, diventa ministro della pubblica istruzione, fa riforma gentile nell'istruzione \t nonostante delitto matteotti lui continua ad essere fedele al fascismo (anche se alcuni dicono che prende le distanza, anche se continua ad essere fascista) \\
Nel 25 diventa direttore dell'istituto Treccani \t elabora e pubblica l'enciclopedia treccani \\
Nel 43 lui aderisce alla repubblica di Salò \t fa atto estremo di un regime di cui era diventato il leader culturale \t lui definisce dal punto di vista ideologico il fascismo e aveva un'egemonia culturale \\
Viene ucciso nel 44 da uno sconosciuto davanti alla sua casa di Firenze
\end{document}
